% Portada de C3: commits, una rama que se une de nuevo y una etiqueta de versión.
\documentclass[tikz,border=0pt]{standalone}
\usepackage{resonancia}
\begin{document}
\begin{tikzpicture}
  \clip[rounded corners=6pt] (0,0) rectangle (12,3.6);
  \shade[top color=rGris!22, bottom color=white] (0,0) rectangle (12,3.6);
  % línea principal
  \draw[trazo, color=rAzul] (1,1.3) -- (11,1.3);
  \foreach \x in {1,3,5,7,9,11} { \fill[white, draw=rAzul, line width=1.1pt] (\x,1.3) circle (0.17); }
  % rama
  \draw[trazo, color=rVerde] (3,1.3) .. controls (3.6,2.5) and (4.2,2.5) .. (4.8,2.5) -- (6.2,2.5) .. controls (6.8,2.5) and (7.4,2.5) .. (7,1.3);
  \foreach \x in {4.8,6.2} { \fill[white, draw=rVerde, line width=1.1pt] (\x,2.5) circle (0.17); }
  \node[etiqueta, color=rVerde, anchor=south] at (5.5,2.75) {feature/pico-a-pico};
  % etiqueta de versión
  \node[caja, draw=rNaranja, fill=rNaranja!10, inner sep=3pt, font=\sffamily\small] at (11,2.2) {v0.1.0};
  \draw[trazo, color=rNaranja] (11,1.5) -- (11,1.95);
  \node[etiqueta, color=rAzul, anchor=north] at (1,1.0) {main};
  \node[etiqueta, color=rGris, anchor=north] at (6,0.55) {historia: cada punto es un commit};
\end{tikzpicture}
\end{document}
